\section{Thinking in Dual Alphabets}\label{sec:dual}
  \begin{table}[b!]
    \centering
    %
    \caption{
      \label{tab:dual-conventions}
      Our dual-alphabet convention, and how to use it.
    }
    %
    \rowcolors{2}{}{gray!10}
    {\setlength{\tabcolsep}{6pt}
    \begin{tabularx}{\textwidth}{XXX}
      \toprule
      \textbf{Convention} & \textbf{Rule/definition} & \textbf{Notes/examples} \\
      \midrule
      Hat $=$ alphabet flag 
        & $\hat{x}$ is the same value as $x$, expressed as a 
          bitstring over the dual alphabet $\{+,-\} = \hbin$.
        & Example: If $x = 5 = \bits{101}$, then 
          $\ket{\hat{x}} = \ket{\hat{5}}
          = \ket{\hone \hnil \hone} = \ket{-+-}\, .$ \\
      Dual-domain states 
        & $\ket{\hat{b}} \coloneqq \frac{1}{\sqrt 2} (\ket{\nil} + (-1)^b \ket{\one})$,
          $\ket{\hat{x}} = \Hgate^{\otimes n} \ket{x}$.
        & Related to how the QKD literature writes $\Hgate^\theta \ket x$ for a state
          expressed in the $\theta$ basis. \\
      Hats are state labels
        & Hatted values never appear in arithmetic outside of kets.
        & Forbidden: $(-1)^{\hat{x} \cdot y}$, $f(\hat{x})$. \\
      Wide hats go over arithmetic expressions.
        & Hats are applied \emph{last}.
        & Examples: $\ket{\widehat{f(x)}}$, $\ket{\widehat{x \xor y}}$. \\
      Xor of labels
        & $\hat{x} \xor \hat{y} \coloneqq \widehat{x \xor y}$.
        & The one exception to the above ``no
          arithmetic'' rule; $\widehat{(\cdot)}$ is a group homomorphism. \\
      Fresh letters for dummy indices
        & Write $\ket{x} = \nicefrac{1}{\sqrt{2^n}} \sum_z (-1)^{x \cdot z} \ket{\hat z}$.
        & $\ket{x} = \nicefrac{1}{\sqrt{2^n}} \sum_{\hat{x}} (-1)^{x \cdot \hat{x}} \ket{\hat x}$ 
          would not be well-defined under our convention. \\
      Outcomes are plain
        & Measurement results always live in $\bin^n$.
        & Write $x \hobstore \register[X]$, not $\hat x \hobstore \register[X]$. \\
      \bottomrule
    \end{tabularx}}
  \end{table}

  \subsection{Dual Alphabets}
    We introduce the dual-alphabet notation, which allows us to reason about qubit
      states in both the computational and the Hadamard basis without having to
      explicitly apply the Hadamard transformation every time.
      Recall the the Hadamard basis states: 
      \[
        \ket{+} \coloneqq \frac{\ket{\nil} + \ket{\one}}{\sqrt{2}}\, ,\quad
        \ket{-} \coloneqq \frac{\ket{\nil} - \ket{\one}}{\sqrt{2}}\, .
      \]
      For $b \in \bin$, define the state
      \[
        \ket{\hat{b}} \coloneqq \frac{\ket{\nil} + (-1)^b \ket{\one}}{\sqrt 2}\, ,
      \]
      i.e.\ $\ket{\hat{b}} = \ket +$ if $b=\nil$, and $\ket{\hat{b}} = \ket -$ if
      $b=\one$. Notice how we are invited to think of the variable $\hat{b}$ as a bit
      defined over the alphabet $\hbin \coloneqq \{+,-\}$ in place
      of our usual alphabet $\bin$, since $\ket{\hnil} = \ket +$ and 
      $\ket{\hone} = \ket -$. We will refer to this as the \emph{dual} alphabet.
      And just as $\ket{\hat{b}}$ can be expressed in the usual alphabet $\bin$
      using the definition above, so can a state $\ket b$ be expressed in terms
      of the dual alphabet $\hbin$:
      \[
        \ket{\hat{b}} = \frac{\ket{\nil} + (-1)^b \ket{\one}}{\sqrt 2}\, ,
        \quad
        \ket{b} = \frac{\ket{\hnil} + (-1)^b \ket{\hone}}{\sqrt 2}\, .
      \]

    Going to multiple qubits, Hadamard-transforming an
      $n$-qubit computational-basis state $\ket{x} = \ket{\bits{101} \ldots \bits{01}}$ just 
      applies the Hadamard gate to each qubit individually:
    \[
      \ket{x} = \ket{\bits{101} \ldots \bits{01}}
        \apply{\Hgate^{\otimes n}}\ket{- + - \ldots + -} 
        = \ket{\hone \hnil \hone \ldots \hnil \hone} 
        \eqqcolon \ket{\hat{x}}\, .
    \]
    The result is a state labelled by a \emph{bitstring} $\hat{x}$, 
      again defined over the dual alphabet: $\hat{x} \in \hbin^n = \{+,-\}^n$.
      Hence, our convention generalizes straightforwardly:
      Writing out the state $\ket{\hat{x}}$ in
      terms of the computational basis, we see that we can \emph{translate} a
      state between alphabets using the usual rule for applying a
      Hadamard-transformation, expressed in the below definition.
    \begin{definition}\label{def:HadamardTranslation}
      Let $x \in \bin^n$. Then,
      \[
        \ket{\hat{x}} = \frac{1}{\sqrt{2^n}} \sum_{z \in \bin^n} (-1)^{x \cdot
        z} \ket{z}\, ,
        \quad 
        \ket{x} = \frac{1}{\sqrt{2^n}} \sum_{z \in \bin^n} (-1)^{x \cdot
        z} \ket{\hat{z}}\, .
      \]
    \end{definition}
    Note how hats only ever show up inside kets: they are to be used as
      \emph{state labels}, never in the surrounding arithmetic.
      Relatedly, since the operation $\widehat{(\cdot)}$ takes any value
      expressed over $\bin$ and translates it to that same value expressed over
      $\hbin$, dummy variables must get a new letter $z$. Finally, classical
      values such as observed outcomes are always defined over $\bin$ and may never carry hats. 
      See \cref{tab:dual-conventions} for a helpful
      overview of how to use our conventions.

    To re-iterate the point: We are not \emph{transforming} the states in the
      above identities
      %\emph{translating} a state is not the same as \emph{transforming} it 
      -- not a single qubit gate was applied -- 
      we are merely looking at the same state with a new pair of glasses!
      This is what allows us to express these as \emph{identities},
      rather than as an \emph{update} (with arrows in place of equality signs).

    So, measuring a state $\ket x$ in the computational basis yields
      outcome $x$ with certainty, measuring a state $\ket{\hat{x}}$ in the
      computational basis yields a uniformly random outcome $x^\prime \sample \bin^n$;
      and while measuring a state $\ket{\hat{x}}$ in the Hadamard basis yields
      outcome $x$ with certainty, measuring a computational-basis state $\ket{x}$ in the
      Hadamard basis will \emph{also} yield a uniformly random outcome $x^\prime \in \bin^n$.

    Over the course of this section, we hope to convince you that
      the ability to translate between alphabets on the fly, expressing any
      given state in whichever basis it is the simplest (even at times using
      different alphabets for different registers), will turn out to be a
      very useful trick that simplifies many intermediate calculations, 
      while bringing to the forefront many of the unique aspects that separate
      the quantum and classical worlds.
      Indeed, experience tells us that
      playing around with the dual alphabets can be very 
      helpful in building that coveted quantum intuition! 

  \subsection{Dual-Alphabet Pseudocode}
    We have introduced the notation $\register[B] \xstore \register[A]$ for the
      operation that, if register $\register[A]$ is in state $\ket a$ and
      register $\register[B]$ is in state $\ket b$, applies the operation 
      \[
        \ket{a,b}_{\register[AB]} \to \ket{a, b \xor a}_{\register[AB]}\, .
      \]
      If $a$ and $b$ are bits ($a,b \in \bin$), then this is just a CNOT
      from qubit $\register[A]$ to qubit $\register[B]$.
      For any function $f$ that we know how to compute, 
      writing $\register[B] \xstore f(\register[A])$ denotes the standard 
      reversible function evaluation:
      \[
        \ket{a,b}_{\register[AB]} \to \ket{a, b \xor f(a)}_{\register[AB]}\, .
      \]

    What are the equivalent operations for Hadamard-basis states? Well, we know
      that the $\Zgate$ (phase-flip) gate plays the role of the $\Xgate$ (not)
      gate in the Hadamard basis; in fact, with the bases being dual, $\Zgate$
      and $\Xgate$ switch roles exactly when moving between the computational
      and Hadamard bases:
      \begin{align*}
        \ket{\nil} \apply{\Xgate} \ket{\one}\, ,\ \ket{\one} \apply{\Xgate} \ket{\nil} 
          &\qquad \ket{\nil} \apply{\Zgate} \ket{\nil}\, ,\ \ket{\one} \apply{\Zgate} -\ket{\one} \\
        \ket{+} \apply{\Zgate} \ket{-}\, ,\ \ket{-} \apply{\Zgate} \ket{+}
          &\qquad \ket{+} \apply{\Xgate} \ket{+}\, ,\ \ket{-} \apply{\Xgate} -\ket{-}
      \end{align*}
      which in our notation may be written more succinctly as
      \begin{align*}
        \ket{b} \apply{\Xgate} \ket{b \xor \one}\, &,\quad \ket{b} \apply{\Zgate} (-1)^b \ket{b} \\
        \ket{\hat{b}} \apply{\Zgate} \ket{\hat{b} \xor \hone}\, &,\quad \ket{\hat{b}} \apply{\Xgate} (-1)^b \ket{\hat{b}}
      \end{align*}
      We want to identify the equivalent of the CNOT operation,
      a.k.a.\ $\CXgate$, in the Hadamard domain, i.e.\ the operation
      \[
        \ket{\hat{a}, \hat{b}} \to \ket{\hat{a}, \hat{b} \xor \hat{a}}\, .
      \]
      This \emph{almost} looks like a controlled-$\Zgate$, or $\CZgate$, except that we
      are controlling on $\ket{\hat a} = \ket{\hone}$. In other words, the
      $\Zgate$ operation is only applied in the branches of the superposition
      where register $\register[A]$ finds itself in state $\ket{\hone} = \ket{-}$ state. 
      Call this operation $\hCZgate$, defined as:
      \[
        \hCZgate_{\register[AB]} = 
          \ketbra{+}_{\register[A]} \otimes \unitary[I]_{\register[B]}
          + \ketbra{-}_{\register[A]} \otimes \Zgate_{\register[B]}\, ,
      \]
      i.e.\ the unitary that picks out the $\ket{-}_{\register[A]}$ branches and
        applies $\Zgate$ to register $\register[B]$ there, and elsewhere applies the
        identity. To see that this is the correct operation, simply notice that since 
        $\Xgate$ and $\Zgate$ switch roles between domains, 
        we might as well write $\Zgate = \hat{\Xgate}$, in which case
        $\hCZgate = \hat{\unitary[C]}\hat\Xgate$.

      Still, it is instructive to see explicitly how it all works out, so let us apply
        $\hCZgate$ to the state $\ket{\hat a, \hat b}$ and see what happens
        by relying on the standard phase-flipping effect of $\Zgate$ on
        computational-basis states:
      \begin{multline*}
        \ket{\hat a, \hat b} 
          = \ket{\hat a} \left( \frac{\ket{\nil} + (-1)^b \ket{\one}}{\sqrt 2} \right) \\
          \apply{\hCZgate} \ket{\hat a} \left( \frac{\ket{\nil} + (-1)^b(-1)^a \ket{\one}}{\sqrt 2} \right) 
          = \ket{\hat a} \left( \frac{\ket{\nil} + (-1)^{b \xor a} \ket{\one}}{\sqrt 2} \right)
          = \ket{\hat a, \widehat{b \xor a}} \eqqcolon \ket{\hat a, \hat{b} \xor \hat{a}}\, ,
      \end{multline*}
      where we have defined the xor between two dual-alphabet strings 
        $\hat a \xor \hat b$ to simply be the dual-alphabet expression of the
        xored value $\widehat{a \xor b}$. (In other words,
        $\widehat{(\cdot)}$ is
        a group homomorphism over $(\ZZ_2, \xor)$.) By translating back to the dual
        alphabet in the last step, we find that this is indeed the operation we
        sought, namely the dual-domain equivalent of the CNOT (or $\CXgate$) operation. 

      For multi-qubit states, recall that the operation 
        $\outReg \xstore \inReg$ on two equally-sized registers $\inReg$ and
        $\outReg$, denotes the operation:
        \[
          \ket{x,y} \to \ket{x,y \xor x}\, ,
        \]
        where now, $y \xor x$ denotes the bitwise xor of $x$ and $y$.
        %(or equivalently, the sum $x + y \pmod{2}$, if we choose to
        %interpret $x$ and $y$ as integers over $\ZZ_{2^n}$ instead of
        %as fixed-length bitstrings over $\bin^n$).
        Hence, this operation can be implemented by simply applying $\CXgate$
        gates between corresponding qubits, so that $\ket{x_i,y_i} \to \ket{x_i,
        y_i \xor x_i}$. It is clear that the corresponding
        Hadamard-domain operation must be to apply $\hCZgate$ between
        the corresponding qubits of $\inReg$ and $\outReg$, so that 
        $\ket{\hat{x}_i, \hat{y}_i} \to \ket{\hat{x}_i, \hat{y}_i \xor
        \hat{x}_i}$, and
        \[
          \ket{\hat x, \hat y} \to \ket{\hat x, \widehat{x \xor y}} 
          \eqqcolon \ket{\hat x, \hat{x} \xor \hat{y}}\, . 
        \]
        We will denote this operation $\outReg \zstore \inReg$, where the $\ominus$
        indicates the $\Zgate$ operation (in reference to its association with
        phase-flipping), and the hat serves as a reminder that we are
        controlling on the states $\ket\hone$ rather than $\ket\one$.\footnote{
          This suggests the additional operators $\czstore$ and $\hxstore$, which
          would indeed be perfectly well-defined, though we will not need them.
        }

  \subsection{Properties of $\xstore$ and $\zstore$.}
    Let us explore how $\xstore$ and $\zstore$ act on various initial states,
      starting with single qubits.
      By design, the $\xstore$ operator takes states $\ket{a, b}$ to
      $\ket{a, b \xor a}$, while 
      $\zstore$ takes states $\ket{\hat a, \hat b}$ to 
      $\ket{\hat a, \hat b \xor \hat a}$. 
      What about the other way around?
      Or for that matter, mixed-alphabet states like $\ket{\hat a, b}$ and
      $\ket{a, \hat b}$? 
      In other words, we would like to fill out the missing
      cells in the following table:

    \begin{steptable}
      \setlength{\extrarowheight}{2pt}
      \rowcolors{2}{}{gray!10} % zebra striping
      \begin{tabular}{ccc}
          \toprule
          $\register[AB]$ initial state & 
            $\quad$ $\register[B] \xstore \register[A]$ yields $\quad$ & 
            $\register[B] \zstore \register[A]$ yields\\
          \midrule
          $\ket{a,b}$ & $\ket{a,b \xor a}$ & \\
          $\ket{a,\hat b}$ &  &  \\
          $\ket{\hat a, b}$ &  &  \\
          $\ket{\hat a, \hat b}$ &  & $\ket{\hat a, \hat b \xor \hat a}$ \\
          \bottomrule
        \end{tabular}
    \end{steptable}

    \subsubsection{Case 1: $\ket{a, b}$, $\register[B] \zstore \register[A]$.}
      Let us begin by asking what is perhaps the most pressing question: How
        does $\zstore$ affect computational-basis 
        states like $\ket{a,b}$? 
        %$\register[A]$ is a single-qubit register 
        %initialized to $\ket a$, and $\register[B]$ is a single-qubit
        %register initialized to $\ket b$, and $a,b \in \bin$. 
        Recall that $\zstore$ does not control on $\ket a = \ket \one$, but
        rather on $\ket{\hat a} = \ket{\hat 1}$, so this does
        \emph{not} act like a standard controlled-$\Zgate$ operation.
        To reveal the relevant controls, begin by 
        translating $\ket a$ to the dual alphabet:
      \[
        \ket{a,b} = \frac{\ket{\hnil} + (-1)^a \ket{\hone}}{\sqrt 2}\ket b\, .
      \]
      For single-qubit registers, the operation $\register[B] \zstore
        \register[A]$ just applies the $\hCZgate$ operation we defined earlier,
        which only fires in the $\ket \hone$ branch of the superposition.
        Meanwhile, $\Zgate$ applied to a computational-basis state $\ket
        b$ flips the phase iff $b = \one$. The result:
      \[
        \frac{\ket{\hnil} + (-1)^a \ket{\hone}}{\sqrt 2}\ket b
          \apply{\hCZgate} 
          \frac{\ket{\hnil} + (-1)^{a} (-1)^{b} \ket{\hone}}{\sqrt 2}\ket b
          = \frac{\ket{\hnil} + (-1)^{a \xor b} \ket{\hone}}{\sqrt 2}\ket b
          = \ket{a \xor b, b}\, ,
      \]
      where we used the translation rule a second time to translate the state back to the
        computational-basis.

      So, the operation $\register[B] \zstore \register[A]$ takes a state
        $\ket{a,b}$ to the state $\ket{a \xor b, b}$. In other words, it acts
        exactly like the operation $\register[A] \xstore \register[B]$, our
        previous operation with the registers flipped! 
        Since operators $\zstore$ and $\xstore$ switch
        roles when moving from the computational basis (standard alphabet)
        to the Hadamard basis (dual alphabet), we may conclude right away 
        that the operation $\register[B] \xstore \register[A]$ must similarly
        take a state $\ket{\hat a, \hat b}$ to the state $\ket{\hat a \xor \hat b, \hat b}$,
        i.e.\ the same as $\register[A] \zstore \register[B]$.
        Let's add both to our table:
      \begin{steptable}
        \setlength{\extrarowheight}{2pt}
        \rowcolors{2}{}{gray!10} % zebra striping
        \begin{tabular}{ccc}
          \toprule
          $\register[AB]$ initial state & 
            $\quad$ $\register[B] \xstore \register[A]$ yields $\quad$ & 
            $\register[B] \zstore \register[A]$ yields\\
          \midrule
          $\ket{a,b}$ & $\ket{a,b \xor a}$ & $\ket{a \xor b, b}$ \\
          $\ket{a,\hat b}$ &  &  \\
          $\ket{\hat a, b}$ &  &  \\
          $\ket{\hat a, \hat b}$ & $\ket{\hat a, \hat b \xor \hat a}$ & $\ket{\hat a, \hat b \xor \hat a}$ \\
          \bottomrule
        \end{tabular}
      \end{steptable}

      The table surfaces two symmetries: First, the diagonal hat/non-hat symmetry
        from swapping $a,b,\xstore \leftrightarrow \hat a, \hat b, \zstore$, 
        which is simply saying that the computational basis and the
        Hadamard basis are dual, and that $\xstore$ and $\zstore$ play the role
        of the other in the other domain.

      The more interesting symmetry is the one we just discovered, that switching $\xstore$ for
        $\zstore$ \emph{and} reversing direction leaves the
        operation invariant. 
        %In other words, applying $\zstore$ from $\inReg$ to
        %$\outReg$ \emph{has the same effect as} applying $\xstore$ from
        %$\outReg$ to $\inReg$.
        Since the symmetry holds for every state in the qubit
        basis $\{\ket{\nil}, \ket{\one}\}$, the symmetry must hold for \emph{every}
        qubit state; and since the $\xstore$ and $\zstore$ operations are
        implemented bitwise for multi-qubit registers, the same must be true in
        general. %Let us state this as a lemma.

      \begin{lemma}\label{lemma:xzstore}
        Let $\inReg$ and $\outReg$ be $n$-qubit. Then
        applying $\xstore$ from $\inReg$ to $\outReg$ is equivalent to applying
        $\zstore$ from $\outReg$ to $\inReg$, and vice versa:
        \[
        \outReg \xstore \inReg\ \equiv\ \inReg \zstore \outReg\, ,
        \quad
        \inReg \xstore \outReg\ \equiv\ \outReg \zstore \inReg\, .
        \]
      \end{lemma}

      \cref{lemma:xzstore} may seem surprising, but it is
        actually just a restatement of the well-known fact that
        CNOT switches direction under the Hadamard transformation -- the target
        becomes the control, and the control becomes the target:
      $
        \Hgate^{\otimes 2}\, \CXgate_{\register[AB]}\, \Hgate^{\otimes 2} = \CXgate_{\register[BA]}\, .
      $

    \subsubsection{Case 2: $\ket{a, \hat b}$, $\register[B] \xstore \register[A]$.}
      Next up is our first multi-alphabet state,
        $\ket{a, \hat b}$. What is the effect of the operation 
        $\register[B] \xstore \register[A]$ on this state, 
        i.e.\ a CNOT from $\ket a$ to $\ket{\hat b}$?
        Well, recall that $\Xgate$ and $\Zgate$ play opposite roles in the
        two domains, so applying an $\Xgate$ gate to a state $\ket{\hat b}$
        will have the same effect on it as if we had applied a $\Zgate$ gate
        to $\ket{b}$: If $\hat b = \hone$, then the
        phase flips, and otherwise nothing happens.
        Since this operation is controlled on $a = \one$, we need that
        \emph{both} $a$ and $b$ are $\one$ for the phase to flip; in other
        words, we need that their \emph{product} is $1$:
      \[
        \ket{a, \hat b} \apply{\CXgate} (-1)^{a \cdot b}\ket{a,\hat b}\, .
      \]
      The only possible effect is a flip of the global phase.
        With global phases being unphysical, this means that
        in isolation, the operation acts on the state like the
        identity.\footnote{
          Another way to say this is that $\ket{\hat a, b}$ is an
          eigenstate of the $\hCZgate$ operator, which follows
          from the fact that $\ket b$ is an eigenstate of the $\Zgate$
          operator.
        }
        But of course, we will often \emph{not} be applying these operations in isolation, and
        the moment the operation is controlled on other qubit states (such as
        when applied inside a ``branch loop''), the phase
        is no longer global, and may certainly \emph{not} be ignored.

      Using the diagonal symmetry of our table, we may again immediately
        conclude that the operation $\register[B] \zstore \register[A]$ with
        initial state $\ket{\hat a, b}$ will have the same overall effect:
      \[
        \ket{\hat a, b} \apply{\hCZgate} (-1)^{a \cdot b}\ket{\hat a, b}\, .
      \]
      Let's add both to the table:
      \begin{steptable}
        \centering
        \setlength{\extrarowheight}{2pt}
        \rowcolors{2}{}{gray!10} % zebra striping
        \begin{tabular}{ccc}
          \toprule
          $\register[AB]$ initial state & 
            $\quad$ $\register[B] \xstore \register[A]$ yields $\quad$ & 
            $\register[B] \zstore \register[A]$ yields\\
          \midrule
          $\ket{a,b}$ & $\ket{a,b \xor a}$ & $\ket{a \xor b, b}$ \\
          $\ket{a,\hat b}$ & $(-1)^{a \cdot b} \ket{a, \hat b}$ &  \\
          $\ket{\hat a, b}$ &  & $(-1)^{a \cdot b} \ket{\hat a, b}$ \\
          $\ket{\hat a, \hat b}$ & $\ket{\hat a, \hat b \xor \hat a}$ & $\ket{\hat a, \hat b \xor \hat a}$ \\
          \bottomrule
        \end{tabular}
      \end{steptable}

    \subsubsection{Case 3: $\ket{\hat a, b}$, $\register[B] \xstore \register[A]$.}
      Finally, what is the effect of applying $\register[B] \xstore
        \register[A]$ on the initial state $\ket{\hat a, b}$? 
        This one \emph{should} look familiar to you: We
        are applying a CNOT from a qubit in superposition to a qubit in a
        computational basis state -- which is exactly how entangled qubit pairs are
        usually prepared, just with $a=\nil$, $b=\nil$!
        In that case, we get the familiar maximally-entangled ``Bell state'' 
        $\nicefrac{1}{\sqrt 2}\, (\ket{\bits{00}} + \ket{\bits{11}})$.
        Let's look at the general case:
        register $\register[A]$ is initialized to $\ket{\hat a}$ and register
        $\register[B]$ is initialized to $\ket b$. Since $\xstore$ controls on
        computational-basis values, begin by expressing $\ket{\hat a}$ in the
        standard alphabet. We get:
        \[
          \ket{\hat a, b} 
            = \frac{\ket{\nil} + (-1)^a \ket{\one}}{\sqrt 2}\ket b\ 
            \apply{\CXgate}\ \frac{\ket{\nil, b} + (-1)^a \ket{\one, b \xor \one}}{\sqrt 2}\, .
        \]
        The result is one of the four states  of the \emph{Bell basis}, 
          %$\{ \ket{\mathrm{\Phi}^+}, \ket{\mathrm{\Phi}^-}, \ket{\mathrm{\Psi}^+}, \ket{\mathrm{\Psi}^-} \}$, 
          depending on the values of $a$ and $b$:
        \begin{align*}
          (a,b) = (\nil,\nil)\ \implies\ \frac{\ket{\bits{00}} + \ket{\bits{11}}}{\sqrt 2} \eqqcolon \ket{\beta_{00}}\, ,
          &\qquad
          (a,b) = (\one,\nil)\ \implies\ \frac{\ket{\bits{00}} - \ket{\bits{11}}}{\sqrt 2} \eqqcolon \ket{\beta_{10}}\, , 
          \\
          (a,b) = (\nil,\one)\ \implies\ \frac{\ket{\bits{01}} + \ket{\bits{10}}}{\sqrt 2} \eqqcolon \ket{\beta_{01}}\, ,
          &\qquad
          (a,b) = (\one,\one)\ \implies\ \frac{\ket{\bits{01}} - \ket{\bits{10}}}{\sqrt 2} \eqqcolon \ket{\beta_{11}}\, .
        \end{align*}
      %Let's adopt the notation of Nielsen and Chuang~\cite{NieChu10} and write
      So, applying $\xstore$ to $\ket{\hat a, b}$ gives the Bell state
        $\ket{\beta_{ab}}$. Diagonal table-symmetry then tells us that applying
        $\zstore$ to the state $\ket{a, \hat b}$
        must yield the same Bell state, but with hats:
        \[
          \ket{a, \hat b} 
            = \frac{\ket{\hnil} + (-1)^a \ket{\hone}}{\sqrt 2}\ket{\hat b}\ 
            \apply{\hCZgate}\ \frac{\ket{\hnil, \hat b} + (-1)^a \ket{\hone, \hat b \xor \hone}}{\sqrt 2}
            \eqqcolon \ket{\hat{\beta}_{ab}}\, .
        \]
        %where we have defined $\ket{\hat{\beta}_{ab}}$ to be the same state with hats.
        %With these added our table is full:
        %\begin{steptable}
        %  \centering
        %  \setlength{\extrarowheight}{2pt}
        %  \rowcolors{2}{}{gray!10} % zebra striping
        %  \begin{tabular}{ccc}
        %    \toprule
        %    $\register[AB]$ initial state & 
        %      $\quad$ $\register[B] \xstore \register[A]$ yields $\quad$ & 
        %      $\register[B] \zstore \register[A]$ yields\\
        %    \midrule
        %    $\ket{a,b}$ & $\ket{a,b \xor a}$ & $\ket{a \xor b, b}$ \\
        %    $\ket{a,\hat b}$ & $(-1)^{a \cdot b} \ket{a, \hat b}$ & $\ket{\hat{\beta}_{ab}}$ \\
        %    $\ket{\hat a, b}$ & $\ket{\beta_{ab}}$ & $(-1)^{a \cdot b} \ket{\hat a, b}$ \\
        %    $\ket{\hat a, \hat b}$ & $\ket{\hat a, \hat b \xor \hat a}$ & $\ket{\hat a, \hat b \xor \hat a}$ \\
        %    \bottomrule
        %  \end{tabular}
        %\end{steptable}
        %
      Interestingly, \cref{lemma:xzstore} tells us that the state
        $\ket{\hat{\beta}_{ab}}$ must \emph{also} be what we'd get if we
        had applied the $\xstore$ operation in reverse, 
        $\register[A] \xstore \register[B]$,
        starting from the same initial state $\ket{a, \hat b}$.
        But clearly, this would just give us
        the same Bell state as before, except with the roles of $a$ and $b$ swapped:
        \[
          \ket{a, \hat b}\ \apply{\CXgate_{\register[BA]}}\
            = \frac{\ket{a, \nil} + (-1)^b \ket{a \xor \one, \one}}{\sqrt 2} =
            (-1)^{a \cdot b}\ket{\beta_{ba}}\, ,
        \]
        where the extra minus sign in the $a = b = 1$ case accounts for the reversed order.
        We are led to conclude that these two states must in fact be \emph{the
        same state}:
        \[
          \ket{\hat{\beta}_{ab}} = \frac{\ket{\hnil, \hat b} + (-1)^a \ket{\hone, \hat b \xor \hone}}{\sqrt 2}
          = \frac{\ket{a, \nil} + (-1)^b \ket{a \xor \one, \one}}{\sqrt 2} =
          (-1)^{a \cdot b}\ket{\beta_{ba}}\, .
        \]
        What we have just discovered is another well-known fact, namely that the Hadamard
          transformation acts like a permutation on the set of Bell states -- and
          that without ever touching a double sum! 
        %\[
        %  \ket{\mathrm{\Phi}^+} \lrapply{\Hgate^{\otimes 2}} \ket{\mathrm{\Phi}^+}\, ,
        %  \quad
        %  \ket{\mathrm{\Phi}^-} \lrapply{\Hgate^{\otimes 2}} \ket{\mathrm{\Psi}^+}\, ,
        %  \quad
        %  \ket{\mathrm{\Psi}^-} \lrapply{\Hgate^{\otimes 2}} \ket{\mathrm{\Psi}^-}\, .
        %\]

        With these additions, our \cref{tab:xzstore} is finally complete.

      \begin{table}[h!]
        \caption{
          \label{tab:xzstore}
          The effect of $\xstore$ and $\zstore$ 
          on the four single-qubit, dual-basis states.
        }
        \centering
        \setlength{\extrarowheight}{2pt}
        \rowcolors{2}{}{gray!10} % zebra striping
          \begin{tabular}{ccc}
            \toprule
            $\register[AB]$ initial state & 
              $\quad$ $\register[B] \xstore \register[A]$ yields $\quad$ & 
              $\register[B] \zstore \register[A]$ yields\\
            \midrule
            $\ket{a,b}$ & $\ket{a,b \xor a}$ & $\ket{a \xor b, b}$ \\
            $\ket{a,\hat b}$ & $(-1)^{a \cdot b}\ket{a, \hat b}$ & $\ket{\hat{\beta}_{ab}} = (-1)^{a \cdot b}\ket{\beta_{ba}}$ \\
            $\ket{\hat a, b}$ & $\ket{\beta_{ab}}$ = $(-1)^{a \cdot b}\ket{\hat{\beta}_{ba}}$ & $(-1)^{a \cdot b}\ket{\hat a, b}$ \\
            $\ket{\hat a, \hat b}$ & $\ket{\hat a \xor \hat b, \hat b}$ & $\ket{\hat a, \hat b \xor \hat a}$ \\
            %\midrule
            \bottomrule
          \end{tabular}
      \end{table}

    \subsubsection{The Multi-Qubit Case.}
      Since the operations are applied bitwise, this
        immediately generalizes to multi-qubit states:
      \[
        \outReg \xstore \inReg : \ket{x,y} \to \ket{x, y \xor x}\, ,
        \quad
        \outReg \zstore \inReg : \ket{\hat x, \hat y} \to \ket{\hat x, \hat y \xor \hat x}\, .
      \]
      \cref{tab:multixzstore} provides a summary.
      \hhnote{Todo: Expand.}

    \begin{table}[h!]
      \caption{
        \label{tab:multixzstore}
        The effect of applying the $\xstore$ and $\zstore$ operator
        on the four possible dual-basis initial states.
        Here, $\ket{\beta_{xy}}$ denotes the generalized Bell states 
        $\ket{\beta_{xy}} \coloneqq 2^{-\nicefrac{n}{2}} \sum_z (-1)^{x \cdot z} \ket{z, y \xor z}$.
      }
      \centering
      %\scriptsize
      \rowcolors{2}{}{gray!10} % zebra striping
      %\resizebox{\textwidth}{!}{
        \begin{tabular}{ccc}
          \toprule
          Initial state & $\register[Y] \xstore \register[X]$ yields & $\register[Y] \zstore \register[X]$ yields\\
          \midrule
          $\ket{x,y}$ & $\ket{x,y \xor x}$ & $\ket{x \xor y, y}$ \\
          $\ket{x,\hat y}$ & $(-1)^{x \cdot y}\ket{x, \hat y}$ & $\ket{\hat{\beta}_{xy}} = (-1)^{x \cdot y}\ket{\beta_{yx}}$ \\
          $\ket{\hat x, y}$ & $\ket{\beta_{xy}} = (-1)^{x \cdot y}\ket{\hat{\beta}_{yx}}$ & $(-1)^{x \cdot y}\ket{\hat x, y}$ \\
          $\ket{\hat x,\hat y}$ & $\ket{\hat x \xor \hat y, \hat y}$ & $\ket{\hat x, \hat y \xor \hat x}$ \\
          %\midrule
          % repeat the table with outputs written in the Hadamard basis?
          \bottomrule
        \end{tabular}
      %} % end resizebox
    \end{table}

  \subsubsection{Function Evaluation.}
    We originally introduced $\xstore$ to do reversible function evaluation,
    $\outReg \xstore f(\inReg)$, so it is natural to wonder how such an
    operation behaves for the various initial-state cases;
    see \cref{tab:fxzstore}.
    \hhnote{Todo: Expand.}
    \hhnote{make filling out \cref{tab:fxzstore} an exercise?}

    \begin{table}[h!]
      \caption{
        \label{tab:fxzstore}
        Function evaluation table.
        Notice how setting $f$ to the identity recovers
        \cref{tab:multixzstore}.
      }
      \centering
      %\scriptsize
      \rowcolors{2}{}{gray!10} % zebra striping
      %\resizebox{\textwidth}{!}{
        \begin{tabular}{ccc}
          \toprule
          Initial state & $\register[Y] \xstore f(\register[X])$ yields & $\register[Y] \zstore f(\register[X])$ yields\\
          \midrule
          $\ket{x,y}$ & $\ket{x,y \xor f(x)}$ & 
            $\nicefrac{1}{\sqrt{2^n}} \sum_z (-1)^{x \cdot z \xor f(z) \cdot y}\ket{\hat{z}, y}$ \\
          $\ket{x,\hat y}$ & $(-1)^{y \cdot f(x)}\ket{x, \hat y}$ & 
            $\nicefrac{1}{\sqrt{2^n}} \sum_z (-1)^{x \cdot z}\ket{\hat z, \hat y \xor \widehat{f(z)}}$
            \\
          $\ket{\hat x, y}$ & 
            $\nicefrac{1}{\sqrt{2^n}} \sum_z (-1)^{x \cdot z}\ket{z, y \xor f(z)}$
            & $(-1)^{y \cdot f(x)}\ket{\hat x, y}$ \\
          $\ket{\hat x,\hat y}$ & 
            $\nicefrac{1}{\sqrt{2^n}} \sum_z (-1)^{x \cdot z \xor f(z) \cdot y}\ket{z, \hat y}$
            & $\ket{\hat x, \hat y \xor \widehat{f(x)}}$ \\
          %\midrule
          % repeat the table with outputs written in the Hadamard basis?
          \bottomrule
        \end{tabular}
      %} % end resizebox
    \end{table}

  \subsection{Demonstrations}
    We next provide some examples of quantum algorithms and protocols expressed
      and analyzed using the conventions we have introduced, in order to
      demonstrate their expressiveness and helpfulness.

    \subsubsection{The Deutsch--Jozsa Algorithm.}
      The Deutsch--Jozsa algorithm learns in a single function call whether a function
        $f: \bin^n \to \bin$ is constant or balanced (evaluates to $\nil$ and
        $\one$ on exactly half of the inputs each), under the promise that it
        is either one or the other. The algorithm is a multi-qubit
        generalization of Deutsch's algorithm, which is notable for being
        the first example of a task on which a quantum computer can provably outperform 
        a classical computer, in the sense that any classical algorithm would need
        two function evaluations (or up to $2^{n-1} + 1$ in
        the general case) to arrive at the answer.

      The procedure is given in pseudocode in \cref{fig:DeutschJozsa}, shown in two
        equivalent implementations depending on whether you start in the
        Hadamard basis and apply computational-basis function evaluations (which
        is how the algorithm is usually presented), or if you start from the
        computational basis and apply a Hadamard-basis function evaluation,
        which brings a slightly different perspective to the same procedure.

      \begin{figure}[h!]
   	    \begin{center}
  	  	  \fbox{
  	  	  	\small
  	  	  	\nicoresetlinenr
  	  	  	\begin{minipage}[t]{3cm}
              \underline{Deutsch--Jozsa}
  	  	  		\begin{nicodemus}
                \smallskip
                \item $\register[XB] \gets \ket{\hnil^n, \hone}$
                \smallskip
                \item $\register[B] \xstore f(\register[X])$
                \smallskip
                \item $x \hobstore \register[X]$
                \smallskip
                \item $\pcreturn x$
  	  	  		\end{nicodemus}
  	  	  	\end{minipage}
            \quad 
  	  	  	\nicoresetlinenr
  	  	  	\begin{minipage}[t]{3cm}
              \underline{Alt.\ implementation}
  	  	  		\begin{nicodemus}
                \smallskip
                \item $\register[XB] \gets \ket{\nil^n, \one}$
                \smallskip
                \item $\register[B] \zstore f(\register[X])$
                \smallskip
                \item $x \obstore \register[X]$
                \smallskip
                \item $\pcreturn x$
  	  	  		\end{nicodemus}
  	  	  	\end{minipage}
          }
        \end{center}
        \caption{
          \label{fig:DeutschJozsa}
          The Deutsch-Jozsa algorithm for determining whether a given function
          $f: \bin^n \to \bin$ is constant or balanced.
        }
      \end{figure}

      Let us analyze the algorithm via its computational-basis implementation
        shown on the right.
        %(\cref{fig:DeutschJozsa}, right).
        The operation $\zstore$ applies the $\hCZgate$
        bitwise from the control to the target, conditioned on the control being
        in state $\ket{\hone}$, so let us start by rewriting the
        initial state of the $\register[X]$ in the Hadamard basis. Writing $N$
        for $2^n$ as usual, we have:
      \[
        \ket{\nil^n, \one} 
          %= \ket{0, \bits{1}} 
          %= \frac{1}{\sqrt{2^n}} \sum_{x \in \{\nil^n, \ldots, \one^n\}}
          %  \ket{\hat{x}, \one}
          %= \frac{1}{\sqrt{2^n}} \sum_{x = 0}^{2^n - 1}
          %  \ket{\hat{x}, \one}
          = \frac{1}{\sqrt{2^n}} \sum_{x \in \ZZ_{N}}
            \ket{\hat{x}, \one}\, .
      \]
      %Recall that the operation $\register[B] \zstore f(\register[X])$ is
      %  shorthand for the operation that computes and stores $\ket{f(x)}$ in an
      %  ancilla register, then applies $\zstore$ from that register to the
      %  target register, before uncomputing the ancilla register state and
      %  discarding it. 
      As shown in \cref{tab:fxzstore} (third row, second column), this gives us:
      \[
        \frac{1}{\sqrt{2^n}} \sum_{x \in \ZZ_{N}} 
          \ket{\hat{x}, \one}
        \apply{\hCZgate^{\otimes n}} 
        \frac{1}{\sqrt{2^n}} \sum_{x \in \ZZ_{N}} 
          (-1)^{f(x)} \ket{\hat x, \one}\, .
      \]
      (We could also have used the second column of the first row to arrive at
      this state directly.)
      If $f(x) = c$ for some constant $c$, then we have the state
      \[
        \frac{1}{\sqrt{2^n}} \sum_{x \in \ZZ_{N}} 
          (-1)^{f(x)} \ket{\hat x, \one}
        = \frac{1}{\sqrt{2^n}} \sum_{x \in \ZZ_{N}} 
          (-1)^{c} \ket{\hat x, \one}
        \cong 
        \frac{1}{\sqrt{2^n}} \sum_{x \in \ZZ_{N}} 
          \ket{\hat x, \one}
        = \ket{\nil^n, \one}\, ,
      \]
      where we factored away the global phase. Hence,
      measuring the $\register[X]$ register yields the outcome $\nil^n$
      with probability one.

    To see what happens in the case that $f(x)$ is \emph{balanced}, first
      translate the state back to the computational basis:
      \[
        \frac{1}{\sqrt{2^n}} \sum_{x \in \ZZ_{N}} 
          (-1)^{f(x)} \ket{\hat x, \one}
        = \frac{1}{\sqrt{2^n}} 
          \sum_{z \in \ZZ_{N}} 
          \frac{1}{\sqrt{2^n}} 
          \sum_{x \in \ZZ_{N}}
          (-1)^{f(x) \xor x \cdot z} \ket{z, \one}\, ,
      \]
      flip the order of summation and define the amplitudes $\alpha_z$ for the 
      outcomes of a computational-basis measurement,
      \[
        = \frac{1}{\sqrt{2^n}} 
          \sum_{z \in \ZZ_{N}} 
          \frac{1}{\sqrt{2^n}} 
          \sum_{x \in \ZZ_{N}}
          (-1)^{f(x) \xor x \cdot z} \ket{z, \one}
        \eqqcolon \frac{1}{\sqrt{2^n}} 
          \sum_{z \in \ZZ_{N}} 
          \alpha_z \ket{z, \one}\, .
      \]
      To find the amplitude $\alpha_0$ that the outcome $\nil^n$
      is observed, insert for $z = 0$:
      \[
        \alpha_0 = \frac{1}{\sqrt{2^n}} 
          \sum_{x \in \ZZ_{N}} (-1)^{f(x)} = 0\, .
      \]
      The sum evaluates to zero precisely because $f$ is balanced \emph{and} we
      are summing over an even number of $x$ values, so every term cancels
      against a partner.
      Hence, we are guaranteed to \emph{not} observe the
      outcome $\nil^n$ in the balanced case. When combined with the
      above, this shows that the procedure returns $\nil^n$
      if \emph{and only if} $f$ is constant.

    \subsubsection{Bernstein--Vazirani.}
      Given oracle access to a function $f: \bin^n \to \bin$ such that 
        $f(x) = s \cdot x$ for some secret $s$, 
        the Bernstein--Vazirani algorithm recovers $s$ with a single call to
        the function, where any classical algorithm needs at least $n$ calls. (The
        standard way to solve the problem classically would be to call 
        $f(10 \ldots 0)$ to recover the first bit of $s$, $f(010 \ldots 0)$ to recover
        the second bit, and so on.) 

      The algorithm can be seen as a variant of the
        Deutsch--Jozsa algorithm, and it is notable for having been used to provide the
        first (relativized) separation between the classical complexity class BPP
        and the quantum complexity class BQP~\cite{BerVaz97}.
        The algorithm is given in pseudocode in \cref{fig:BernsteinVazirani}.

      \begin{figure}[h!]
   	    \begin{center}
  	  	  \fbox{
  	  	  	\small
  	  	  	\nicoresetlinenr
  	  	  	\begin{minipage}[t]{3cm}
              \underline{Bernstein--Vazirani}
  	  	  		\begin{nicodemus}
                \smallskip
                \item $\register[XB] \gets \ket{\nil^n, \one}$
                \smallskip
                \item $\register[B] \zstore f(\register[X])$
                \smallskip
                \item $s \obstore \register[X]$
                \smallskip
                \item $\pcreturn s$
  	  	  		\end{nicodemus}
  	  	  	\end{minipage}
            \quad 
  	  	  	\nicoresetlinenr
  	  	  	\begin{minipage}[t]{3cm}
              \underline{Alt.\ implementation}
  	  	  		\begin{nicodemus}
                \smallskip
                \item $\register[XB] \gets \ket{\hnil^n, \hone}$
                \smallskip
                \item $\register[B] \xstore f(\register[X])$
                \smallskip
                \item $s \hobstore \register[X]$
                \smallskip
                \item $\pcreturn s$
  	  	  		\end{nicodemus}
  	  	  	\end{minipage}
          }
        \end{center}
        \caption{
          \label{fig:BernsteinVazirani}
          The Bernstein-Vazirani algorithm for recovering a secret $s$ given
          oracle access to a function $f = s \cdot x$.
        }
      \end{figure}

      Let us analyze it, starting from the implementation on the left.
        The second cell of the first row of \cref{tab:fxzstore} tells us the
        effect of the $\zstore$ operator on a computational basis state;
        inserting for $x = 0$, $y=1$ gives us:
      \[
        \ket{\nil^n, \one}
        \apply{\register[B] \zstore f(\register[X])}
        \frac{1}{\sqrt{2^n}} \sum_{z \in \bin^n} (-1)^{f(z)} \ket{\hat{z}, \one}
        = \frac{1}{\sqrt{2^n}} \sum_{z \in \bin^n} (-1)^{s \cdot z} \ket{\hat{z}, \one}\, ,
      \]
      where we inserted the definition of $f$ in the last step. But looking at
      the state on the right, we can recognize it as simply a
      Hadamard-translated state $\ket{x}$, except that $x = s$:
      \hhnote{Add definition number that I can reference?}
      \[
        \frac{1}{\sqrt{2^n}} \sum_{z \in \bin^n} (-1)^{s \cdot z} \ket{\hat{z}, \one}
        = \ket{s, \one}\, .
      \]
      Applying a computational-basis measurement to the $\register[X]$ register,
      then, immediately yields the secret $s$.

    \subsubsection{Simon's Algorithm.}
      Simon's algorithm is an early example of a quantum period-finding
        algorithm. Given a periodic two-to-one function $f$ with a hidden period
        $s$, so that $f(x) = f(x \xor s)$, it recovers $s$.

      The main quantum step is in recovering a value $x_i$ which is guaranteed to
        satisfy $x_i \cdot s = \nil$. This subroutine is then repeated until enough
        such $x_i$ have been recovered that the equation system is fully determined,
        at which point $s$ may be recovered via basic linear algebra.

      \begin{figure}[h!]
   	    \begin{center}
  	  	  \fbox{
  	  	  	\small
  	  	  	\nicoresetlinenr
  	  	  	\begin{minipage}[t]{3cm}
              \underline{Simon's subroutine}
  	  	  		\begin{nicodemus}
                \smallskip
                \item $\register[XY] \gets \ket{\hnil^n, \nil^n}$
                \smallskip
                \item $\register[Y] \xstore f(\register[X])$
                \smallskip
                \item $x \hobstore \register[X]$
                \smallskip
                \item $\pcreturn x$
  	  	  		\end{nicodemus}
  	  	  	\end{minipage}
            \quad
  	  	  	\nicoresetlinenr
  	  	  	\begin{minipage}[t]{3cm}
              \underline{Alt.\ implementation}
  	  	  		\begin{nicodemus}
                \smallskip
                \item $\register[XY] \gets \ket{\nil^n, \hnil^n}$
                \smallskip
                \item $\register[Y] \zstore f(\register[X])$
                \smallskip
                \item $x \obstore \register[X]$
                \smallskip
                \item $\pcreturn x$
  	  	  		\end{nicodemus}
  	  	  	\end{minipage}
          }
        \end{center}
        \caption{
          \label{fig:Simon}
          The quantum subroutine for Simon's algorithm.
        }
      \end{figure}

      The quantum subroutine is shown in \cref{fig:Simon}. Let us see how it
        gives us an $x$ that is guaranteed to satisfy $x \cdot s = \nil$,
        starting from
        the implementation shown on the left of \cref{fig:Simon}.
      \begin{align*}
        &\ket{\hnil^n, \nil} 
          = \frac{1}{\sqrt{2^n}} \sum_{z \in \ZZ_N} \ket{z, \nil} 
          \apply{\outReg \xstore f(\inReg)} 
          \frac{1}{\sqrt{2^n}} \sum_{z \in \ZZ_N} \ket{z, f(z)}\, .
      \end{align*}
      After this, none of the following operations depend on the state of the $\outReg$ register, 
      so though the procedure in \cref{fig:Simon} is never shown to measure it, it
      makes our analysis a lot easier if we just measure it and discard the
      $\outReg$ register. For each value $y$ that we might observe, there are
      two input values $z$ such that $f(z) = y$, so this is not enough to
      completely collapse the superposition:
      \begin{align*}
        \frac{1}{\sqrt{2^n}} \sum_{z \in \ZZ_N} \ket{z, f(z)}
        &\apply{y \obstore \outReg}
          \frac{
            \ket{z} + \ket{z \xor s}
          }{\sqrt{2}}\, ,
          %\qquad
          %\text{s.th.}\
          %f(z) = y\, .
      \end{align*}
      where $z$ satisfies that $f(z) = f(z \xor s) = y$ for the value $y$ that
      we measured.
      We are building towards a Hadamard-measurement, so translate this state to
      the dual domain so that we can inspect the amplitudes directly:
      \[
          \frac{
            \ket{z} + \ket{z \xor s}
          }{\sqrt{2}}
          = \frac{1}{\sqrt{2^n}} \sum_{x \in \ZZ_N}
          \frac{
            (-1)^{z \cdot x}\ket{\hat x} + (-1)^{(z \xor s) \cdot x}\ket{\hat x}
          }{\sqrt{2}} 
          = \sum_{x \in \ZZ_N}
          \frac{
            (-1)^{z \cdot x}
          }{\sqrt{2^{n+1}}} (1 + (-1)^{x \cdot s}) \ket{\hat x} 
          \eqqcolon \sum_{x \in \ZZ_N} \hat{\alpha}_{x} \ket{\hat x}\, .
      \]
      Look at the amplitude $\hat{\alpha}_x$. Which values of $x$ are we more 
      likely to observe? Well, $\hat{\alpha}_x$ is proportional to the sum $1 +
      (-1)^{x \cdot s}$, where $x \cdot s$ denotes the bitstring inner product.
      There are only two cases: If $x \cdot s = \one$, then $\hat{\alpha}_x = 1-1 = 0$, which means
      that the probability of observing such an outcome is also zero.
      Hence, for any $x$ that we actually do observe, we must have that $x \cdot s =
      \nil$. Indeed, restricting to $x \cdot s = \nil$, we see that we have the state
      \[
        \sum_{x\, :\, x \cdot s = \nil}
        \frac{
          (-1)^{z \cdot x}
        }{\sqrt{2^{n+1}}} (1 + 1) \ket{\hat x}
        = \sum_{x\, :\, x \cdot s = \nil}
        \frac{
          (-1)^{z \cdot x}
        }{\sqrt{2^{n-1}}} \ket{\hat x}\, ,
      \]
      which as a sanity check we can see is normalized, since there are
      $2^{n-1}$ values of $x$ %, i.e.\ half of all $x$, 
      that satisfy $x \cdot s = \nil$. (Otherwise, that would have meant we had
      made an error somewhere.)

    Applying a Hadamard-basis measurement to this state, then, is a procedure to 
      sample uniformly from the set of $x$ that satisfy $x \cdot s = \nil$.
      Importantly, this shows that not only do we get a sample that satisfies
      our requirement, but we are also very unlikely to get the
      \emph{same} $x$ when we later repeat the subroutine, and that
      in all likelihood we will only need to repeat the procedure roughly $n$
      times to get $n$ linearly independent values $x_i$ that all satisfy 
      $x_i \cdot s = \nil$.

    \subsubsection{Grover's Algorithm.}

    \subsubsection{Quantum Teleportation.}

  \clearpage
  \subsection{The Primer Paragraphs}
    \hhnote{
      copy-pasted while I decide what to keep where
    }
    \subsubsection{The Hadamard Basis.}
      Actually, we don't have to imagine explicitly applying the $\Hgate$ gates
        to our qubits before measuring them, as Hadamard-basis measurements
        are perfectly well-defined on their own merits, as one choice among a
        continuous infinity of measurements we could choose to apply.
	      %(we will make this formal in \cref{sec:measurements}).
        Perhaps it is better, then, to view the Hadamard basis as a sort of
        ``dual domain'' under which we can view our qubit,
	      with two possible questions we can ask our state: 
	      Either ``Are you $\ket{\nil}$ or $\ket{\one}$?'', or 
	      ``Are you $\ket +$ or $\ket -$?'' Either way, 
	      the qubit will provide a definite answer, after which its state will 
        ``collapse'' to the answer you got.

	    This is a point worth internalizing, so let us repeat it: 
	      \emph{The set of possible post-measurement states changes to match
	      the question you chose to ask.} If you ask ``Are you $\nil$ or $\one$?'', 
	      then the post-measurement state will be $\ket{\nil}$ or $\ket{\one}$; 
	      if you ask ``Are you $+$ or $-$'', then the post-measurement state will
        be $\ket +$ or $\ket -$. 
        So although at one level, the ``collapse'' of the quantum state is just
        the quantum version of how one would use Bayes' theorem 
        to update a prior probability distribution to a marginal distribution
        conditioned on the outcome, at another level, it represents the
        deepest and strangest mystery in all of quantum physics.
        The grand mystery of quantum physics really comes down to this one fact, 
        that the outcome statistics and the post-observation state both depend
        on \emph{which measurement you chose to apply}, and it has led to a century
        of heated philosophical debate over how quantum theory should be
        interpreted. Ask five physicists for their opinion on the matter, and
        you are likely to get five different answers!\footnote{
          %After we have observed the outcome, 
          %the superposition state ``collapses'' to the corresponding basis state
          %$\ket b$, so that the updated physical system is guaranteed to be consistent with
          %the value we observed. 
          The term ``collapse'' that we
          used above is typically associated with the rather conservative
          \emph{Copenhagen interpretation}, championed by Niels Bohr, and commonly
          characterized as the ``shut-up-and-calculate'' approach, since
          it simply posits that ``such questions are disallowed''.
          Another popular interpretation is the many-worlds interpretation due
          to Hugh Everett, favoured by father of quantum computing theory
          David Deutsch among others. Instead of the quantum state
          ``collapsing'' to a ``classical'' state, this interpretation says that
          both terms are ill-defined, and that the universe should instead be
          understood to be nothing but one big quantum superposition over all
          the worlds that could have possibly arisen from
          the initial conditions of the universe.
          (After all, we the observers are made of quantum particles too, so
          should we not then also be considered part of the quantum state?)
          %; from this point of view, when
          %you made your observation, what \emph{actually} happened
          %was that the quantum state describing the subsystem you
          %were studying and the quantum state describing \emph{you} (after all,
          %you \emph{are} made of quantum particles, are you not?), became
          %completely \emph{entangled}, so that the part of the superposition
          %in which the qubit takes the value $0$ is guaranteed to coincide with
          %the part of the superposition in which you observed the outcome $0$.
          %The world in which the qubit took the value $1$ and you also observed
          %the value $1$ is still ``out there'' in the grand superposition of
          %the universe, it's just been completely cut off from your subjective
          %experience, so from your point of view it might as well not exist.
          If you're curious to learn more about the various interpretations of
          quantum mechanics, and the strengths and weaknesses of each attempt
          at a consistent resolution, Carlo Rovelli's book \emph{Helgoland} provides an
          excellent overview.
          %, and even proposes its own fascinating resolution
          %to the problem.)
          %(In case you're wondering, Hans places his bet on a third option,
          %namely the \emph{relational} interpretation due to Carlo Rovelli,
          %which may be summarized as ``many-worlds without the many worlds''.)
        }
        %(And, of course, the outcome distribution also depends on which question
        %you choose to ask.)

      What about a question like ``Are you $\ket{\nil}$ or
        are you $\ket +$?'' Under our notion of measurement, this question is not
        well-defined, since $\ket{\nil}$ and $\ket +$ are not orthogonal states and
        therefore do not lead to mutually-exclusive outcomes. More on this in
        \cref{sec:measurements}, but for now, think of a measurement as always
        being defined in terms of a set of mutually-exclusive outcomes.

	    Of course, any vector space has a continuous infinity of 
	      bases to choose from, each of which define a binary ``question'' we
	      can pose to our qubit, and the Hadamard
	      basis is just one particularly nice choice:
	      geometrically, it is just the basis rotated by $45^\circ$
	      relative to the original one, as illustrated in
	      \cref{fig:hadamard-basis}.
        %
	      \crakzfigcode{hadamard-basis}{
	        The computational basis $\{\ket{\nil}, \ket{\one}\}$ and the Hadamard basis
	        $\{\ket +, \ket -\}$ describe the same qubit, seen along axes turned
	        by $45^\circ$ relative to each other. 
          %Note how $\ket 0$ bisects
	        %$\ket +$ and $\ket -$ exactly as $\ket +$ bisects $\ket 0$ and
	        %$\ket 1$: the two views are duals of one another.
	      }
        %
	      To appreciate how this makes for a perspective
         ``dual'' to the computational basis,
	      take a look at the following identities:
        \[
          \ket{+} = \frac{\ket{\nil} + \ket{\one}}{\sqrt{2}}\, , \quad 
          \ket{-} = \frac{\ket{\nil} - \ket{\one}}{\sqrt{2}}
          \qquad
          \Longleftrightarrow
          \qquad
          \ket{\nil} = \frac{\ket{+} + \ket{-}}{\sqrt{2}}\, , \quad 
          \ket{\one} = \frac{\ket{+} - \ket{-}}{\sqrt{2}}\, .
        \]
        Notice the symmetry: The computational basis states $\ket{\nil}$ and $\ket{\one}$
        are to the Hadamard basis \emph{exactly} what the plus state and minus
        states are to the computational basis; they really are \emph{dual views},
        which can be translated between via the map $\Hgate$.\footnote{
          The notion of a ``dual domain'' will sound familiar to anyone who has
          ever come across Fourier analysis, and not without reason: The Hadamard
          map is exactly the single-qubit Fourier transformation.
        }
        %It is worth verifying the equivalence by inserting for $\ket +$
        %and $\ket -$ and checking that the extra terms cancel out exactly.

      How does this duality generalize to multi-qubit states?
        We had that
        \[
          \ket{+^n} 
          = \frac{\ket{\nil \ldots \nil\nil} + \ket{\nil \ldots \nil\one} + \ldots + \ket{\one \ldots \one\one}}{\sqrt{2^n}}
          = \sum_{x \in \bin^n} \frac{1}{\sqrt{2^n}}\ket x
        \]
        With $\{\ket{\nil}, \ket{\one}\}$ and $\{\ket +, \ket -\}$ being exact duals of
        each other, then it must also be true that 
        \[
          \ket{\nil^n} 
          = \frac{\ket{+ \ldots ++} + \ket{+ \ldots +-} + \ldots + \ket{- \ldots --}}{\sqrt{2^n}}
          = \sum_{\hat{x} \in \{+,-\}^n} \frac{1}{\sqrt{2^n}} \ket{\hat x}\, .
        \]
        Here, we have introduced a new kind of bitstring, one that is defined
        over the alphabet $\{+, -\}$ rather than the alphabet $\bin$. We will use the
        convention that whenever a number or variable is represented with a hat,
        then it is understood to be a \emph{Hadamard-bitstring} like above. For example,
        whereas we previously represented the decimal $21$ by the bitstring
        $[21]_2 = \bits{10101}$, writing $\widehat{21}$ would imply that it is
        represented by the Hadamard-bitstring $\big[\, \widehat{21}\, \big]_2 = -+-+-$. 

      As we might have anticipated, just like the all-zero state $\ket{\nil^n}$
        maps to the decimal state $\ket 0$, 
        %and the all-ones state $\ket{1^n}$
        %maps to the decimal $(N-1)$-state $\ket{N-1}$ (where $N = 2^n$ as usual),
        the uniform superposition $\ket{+^n}$
        then maps to the Hadamard-zero state $\ket{\hat 0}$.
        %, while the all-minus state
        %$\ket{-^n}$ maps to the state $\ket{\hat{N}-\hat{1}}$.
        Thus, we recover our dual views on the quantum
        state space, which we can translate between as desired via the Hadamard
        transformation. Defining $\hat{\ZZ}_N \coloneqq \{ \hat{x}\ |\ x \in
        \ZZ_N\}$, and letting the dual-alphabet inner bitstring product
        $\hat x \cdot x$ simply treat $\{+,-\}$ as if it were $\bin$,
        we can then express both directions of the Hadamard transformation in
        purely decimal notation:
        \[
          \ket{\hat x} = \sum_{x \in \ZZ_N} 
          \frac{(-1)^{\hat x \cdot x}}{\sqrt{2^n}} \ket x\, ,
          \quad
          \ket{x} = \sum_{\hat x \in \hat{\ZZ}_N} \frac{(-1)^{x \cdot
          \hat x}}{\sqrt{2^n}} \ket{\hat{x}}\, .
        \]
        
      The above notation will at times allow us to greatly simplify our
        notation, making arguments clearer. Consider for example the state
        \[
          \ket{\psi} = 
          -\frac{1}{\sqrt{2^n}} \ket{0}\
          +
          \sum_{x \in \ZZ_N \backslash \{0\}} 
          \frac{1}{\sqrt{2^n}} \ket{x}\, ,
        \]
        i.e.\ the uniform superposition state that singles out the $\ket{0}$ state (implicitly
        understood to be the $\ket{\nil^n}$ state since we are using decimal notation) 
        and applies a negative phase to it. Computing the Hadamard transformation of
        such a state can lead to a mess of double sums, but using the above
        notation, not only is applying the transformation as simple as adding
        hats everywhere,
        \[
          %\ket{\psi} 
          - \frac{1}{\sqrt{2^n}} \ket{0}\
          +
          \sum_{x \in \ZZ_N \backslash \{0\}} 
          \frac{1}{\sqrt{2^n}} \ket{x}
          \quad
          \apply{\Hgate^{\otimes n}} 
          \quad
          - \frac{1}{\sqrt{2^n}} \ket{\hat{0}}\ +
          \sum_{\hat x \in \hat{\ZZ}_N \backslash \{\hat{0}\}} 
          \frac{1}{\sqrt{2^n}} \ket{\hat{x}}\, ,
        \]
        but we can then even choose to write out the state using a mix of
        notations, ``unpacking'' the complex-phase term for closer study while keeping 
        the real-valued terms around in their Hadamard-bitstring form:
        \[
          -\frac{1}{\sqrt{2^n}} \ket{\hat{0}}\ +
          \sum_{\hat x \in \hat{\ZZ}_N \backslash \{\hat{0}\}} 
          \frac{1}{\sqrt{2^n}} \ket{\hat{x}}
          =
          \sum_{x \in \ZZ_N} -\frac{1}{2^n} \ket{x}\ +
          \sum_{\hat x \in \hat{\ZZ}_N \backslash \{\hat{0}\}} 
          \frac{1}{\sqrt{2^n}} \ket{\hat{x}}\, .
        \]

        %\begin{exercise}
        %  Compute the Hadamard transformation of $\ket{\psi}$ using purely
        %  standard bitstring notation, expressed without double sums. You will
        %  need \ldots
        %\end{exercise}
